%%%%%%%%%%%%%%%%%%%%%%%%%%%%%%%%%%%%%%%%%%%%%%%%%%%%%%%%%%%%%%%%%%%%%%
% Lớp 8 | Chương III. Tứ giác | Hình thoi và hình vuông
% Dạng 8H3?14-3: Bài toán tổng hợp về tứ giác đặc biệt
%%%%%%%%%%%%%%%%%%%%%%%%%%%%%%%%%%%%%%%%%%%%%%%%%%%%%%%%%%%%%%%%%%%%%%


% nguồn: De_AI_Toan8_cuoi_hoc_ki_i_1 (nhập 21/11/2026)
\begin{bt}%[8H3H14-3]%[ma:47b4e751]
	Cho tam giác $ABC$ vuông cân tại $A$ có $AB = 8\text{ cm}$. Gọi $M$ là trung điểm của $BC$, lấy điểm $D$ đối xứng với $A$ qua $M$.
a) Chứng minh tứ giác $ABDC$ là hình vuông.
b) Tính độ dài đường chéo $AD$ và diện tích hình vuông $ABDC$.
	\loigiai{
		a) Tứ giác $ABDC$ có hai đường chéo $AD$ và $BC$ cắt nhau tại trung điểm $M$ của mỗi đường nên $ABDC$ là hình bình hành. \\
Vì $\widehat{BAC} = 90^\circ$ nên hình bình hành $ABDC$ là hình chữ nhật. \\
Lại có $AB = AC$ (tam giác $ABC$ vuông cân tại $A$), nên hình chữ nhật $ABDC$ có hai cạnh kề bằng nhau, do đó $ABDC$ là hình vuông. \\
b) Vì $ABDC$ là hình vuông cạnh $8\text{ cm}$ nên $AD = BC = \sqrt{8^2 + 8^2} = 8\sqrt{2}\text{ cm}$. \\
Diện tích hình vuông $ABDC$ là $8^2 = 64\text{ cm}^2$.
	}
\end{bt}

% nguồn: De_AI_Toan8_cuoi_hoc_ki_i_1 (nhập 21/11/2026)
\begin{bt}%[8H3H14-3]%[ma:d0a8238d]
	Cho hình thoi $ABCD$ có hai đường chéo $AC$ và $BD$ cắt nhau tại $O$. Gọi $M$, $N$, $P$, $Q$ lần lượt là trung điểm của các cạnh $AB$, $BC$, $CD$, $DA$.
a) Chứng minh tứ giác $MNPQ$ là hình chữ nhật.
b) Hình thoi $ABCD$ cần thêm điều kiện gì để tứ giác $MNPQ$ trở thành hình vuông?
	\loigiai{
		a) Tương tự bài toán đường trung bình trong các tam giác $ABC$ và $ADC$, ta có $MN \parallel AC$, $MN = \dfrac{1}{2}AC$ và $QP \parallel AC$, $QP = \dfrac{1}{2}AC$, suy ra $MNPQ$ là hình bình hành. \\
Mặt khác $MQ$ là đường trung bình của tam giác $ABD$ nên $MQ \parallel BD$. \\
Vì $ABCD$ là hình thoi nên $AC \perp BD$, suy ra $MN \perp MQ$, tức $\widehat{QMN} = 90^\circ$. \\
Hình bình hành $MNPQ$ có một góc vuông nên $MNPQ$ là hình chữ nhật. \\
b) Để hình chữ nhật $MNPQ$ là hình vuông thì cần $MN = MQ$. \\
Mà $MN = \dfrac{1}{2}AC$ và $MQ = \dfrac{1}{2}BD$, nên điều kiện là $AC = BD$. \\
Hình thoi $ABCD$ có hai đường chéo bằng nhau khi và chỉ khi $ABCD$ là hình vuông.
	}
\end{bt}

% nguồn: chuyen_dang (nhập 22/11/2026)
\begin{ex}%[8H3N14-3]%[ma:14ffef94]
	 Mệnh đề nào sau đây là mệnh đề đúng?
	 \choice
		{Tứ giác có 4 cạnh bằng nhau là  hình vuông}
		{Tứ giác có 2 đường chéo bằng  nhau là hình chữ nhật}
		{\True Hình thang cân có 1  góc vuông là hình chữ nhật}
		{Hình thoi là hình vuông có 1 góc  vuông}
	\loigiai{ 
	 }
\end{ex}

% nguồn: chuyen_dang (nhập 22/11/2026)
\begin{ex}%[8H3H14-3]%[ma:c7176f86]
	 Mệnh đề đảo của mệnh đề nào sau đây là một mệnh đề đúng?
	\choice 
		{Nếu một tứ giác là hình thoi thì tứ giác đó có hai đường  chéo vuông góc}
		{\True Nếu một tứ giác có hai đường chéo cắt nhau tại trung điểm mỗi đường thì tứ giác đó là  hình bình hành}
		{Nếu  một hình thang là hình thang cân thì hình thang đó có hai cạnh bên bằng nhau}
		{Nếu một tứ giác là hình vuông thì tứ giác đó có $4$ cạnh bằng nhau} 
	 \loigiai{
		Nếu một tứ giác có hai đường chéo cắt nhau tại trung điểm mỗi đường thì tứ giác đó là hình bình hành có mệnh đề đảo là  mệnh đề đúng.
	 }
\end{ex}
